% !TeX root = ../../Book3.tex
% 纲要标题：### 12.4 泛自由性与泛平坦性
% 纲要位置：工作记录/计划纲要.md，第 2186 行。
% 纲要细目（写作范围）：
% * 固定正文版本：\(A\) 为 Noether 整环、\(B\) 为有限型 \(A\)-代数、\(M\) 为有限生成 \(B\)-模；证明存在 \(0\ne f\in A\) 使 \(M_f\) 作为 \(A_f\)-模自由，且自由秩未必有限
% * 从 \(\operatorname{Frac}(A)\) 上的有限关系和首项计算出发清除有限个分母；证明中所需的有限自由模项序与 Dickson 型有限性在本节就地建立，不前置第21章的合冲模算法
% * 同时控制代数及其有限模，导出稠密主开集上的平坦性；说明基底整性、有限型与有限生成假设的用途
% * 在已经具有概形语言时再解释一般纤维与特殊纤维；此处只用局部化和张量积表述，不要求先读第24章
% * 代数有限呈示、模有限展示时的非 Noether 推广作为结论范围的补充说明，不与正文版本混写；参考 [V13-3]


\section{泛自由性与泛平坦性}
\label{b3:sec:1204}

设一个方程族的系数属于整环 \(A\)。转到分式域以后，有限关系中的非零系数都可以相除；问题是能否只倒置有限个系数，便在同一个开集上得到一组基。下面用首项与有限分母回答这个问题。证明所需的模单项式在此处建立，不要求预先学习合冲模算法。


\subsection{Noether 整环上的泛自由性定理}
\label{b3:subsec:1204:01}

\begin{theorem}[泛自由性]\label{b3:thm:generic-freeness}
设 \(A\) 为 Noether 整环，\(B\) 为有限型交换 \(A\) 代数，\(M\) 为有限生成 \(B\) 模。则存在 \(0\ne f\in A\)，使 \(M_f\) 作为 \(A_f\) 模自由。其自由秩允许为零或无限。
\end{theorem}

\term[fan-ziyouxing]{泛自由性}[generic freeness]中的“泛”表示可以从基底去掉一个真闭集。这里没有要求 \(M\) 是有限 \(A\) 模；例如 \(B=M=A[x]\) 已经是自由 \(A\) 模，但基 \(1,x,x^2,\ldots\) 为无限集。定理也允许整个模在局部化后消失，例如 \(M=A/(a)\)、\(a\ne0\)，倒置 \(a\) 后自由秩为零。

\ProseParagraphBreak
证明放在下一小节，核心是为商模选出一组可同时保留的单项式。更一般的泛自由性结果可参阅 Matsumura 的《Commutative Ring Theory》\cite[\S24, Theorem 24.1, p. 186]{Matsumura1989CommutativeRingTheory}；这里给出的证明专门处理本章的有限展示数据。

\subsection{有限自由模项序与分母清除}
\label{b3:subsec:1204:02}

令 \(K=\operatorname{Frac}(A)\)，\(P=A[x_1,\ldots,x_s]\)。选满射 \(P\to B\) 及有限组 \(B\) 模生成元，得到 \(M=P^q/N\)。由 \(A\) 的 Noether 性和 Hilbert 基定理，关系模 \(N\) 有限生成。

\begin{lemma}\label{b3:lem:module-leading-finiteness}
设 \(K\) 为域，\(s,q\) 为非负整数，\(e_1,\ldots,e_q\) 为 \(K[x_1,\ldots,x_s]^q\) 的标准基。给模单项式 \(x^\alpha e_j\) 按全次数、指数的字典序、最后按 \(j\) 排序。这是良序，且同乘一个单项式保持大小。任一 \(K[x_1,\ldots,x_s]\) 子模 \(L\) 的非零元素的首单项式，均被有限多个来自 \(L\) 的首单项式之一整除；整除只在同一基向量 \(e_j\) 上比较。
\end{lemma}
\begin{proof}
给定次数的单项式只有有限多个，所以任一非空集合先选最小次数，再选其中最小项，得到良序；加同一个指数向量不改变比较结果。

所需有限性归结为 \(\mathbb N^s\) 中逐坐标偏序的极小元有限。为证明这一点，先注意任一自然数序列都有无限不减子序列：某个值出现无限次时取常值子序列；否则各有限区间只出现有限次，逐次选更大的项。对 \(s\) 个坐标依次取子序列，可知 \(\mathbb N^s\) 的任意无限序列包含两个按原次序排列且逐坐标不减的项。因此一个集合不可能有无限多个不同的极小元。又对任何 \(\alpha\)，其下方的有限盒子 \(\{\beta:\beta\le\alpha\}\) 中可选极小元，故每个元素都在某个极小元之上。这就是所需的 Dickson 型有限性。

分别对每个位置 \(j\) 的首项指数集合应用它，选出有限组极小首项，并取产生它们的 \(g_i\in L\)。全部其他首项都被这些首项之一整除。
\end{proof}

\begin{proof}[\cref{b3:thm:generic-freeness}的证明]
对 \(L=K\otimes_A N\subseteq K[x_1,\ldots,x_s]^q\) 应用引理，选 \(g_1,\ldots,g_t\)，并把首项系数归一成一。逐次消去可被它们的首项整除的最大项，或把不能整除的最大项移入余式。每一步使尚待处理的首单项式严格下降，故有限步终止。对 \(L\) 中的元素，最后余式若非零，其首项一方面属于 \(L\) 的首项，另一方面不被任何所选首项整除，矛盾。因此这些 \(g_i\) 生成 \(L\)。若 \(L=0\)，下述空列表的论证同样成立。

取 \(N\) 的有限组 \(P\) 模生成元 \(v_1,\ldots,v_h\)。每个 \(g_i\) 是 \(v_j\) 的 \(K[x]\) 线性组合，每个 \(v_j\) 也是 \(g_i\) 的 \(K[x]\) 线性组合。这里只涉及有限多个多项式及有限多个系数，故可选非零 \(f\in A\)，倒置它后使所有系数落在 \(A_f\)。于是 \(g_i\in N_f\)，且它们生成 \(N_f\)，首项系数仍为一。

称不被任何 \(g_i\) 的首单项式整除的模单项式为标准单项式。刚才的首项消去只除以首项系数一，故能直接在 \(A_f[x]^q\) 中进行。这说明标准单项式的类生成 \(P_f^q/N_f=M_f\) 作为 \(A_f\) 模。若这些类之间有一个非零有限线性关系，所得非零向量 \(w\in N_f\) 在 \(K[x]^q\) 中仍非零，因为 \(A_f\subseteq K\)。其首单项式是标准的，却又是 \(L\) 中元素的首单项式，与引理矛盾。因此标准单项式的类线性无关，组成所求自由基。
\end{proof}

有限性在此有两个不同用途：Dickson 型论证限制需要的首项数，关系模的有限生成性限制返回原环时需要清除的分母数。只说“清除分母”而不说明哪一组有限数据，不足以得到同一个 \(f\)。

\subsection{稠密主开集上的平坦性}
\label{b3:subsec:1204:03}

\begin{corollary}[泛平坦性]\label{b3:cor:generic-flatness}
设 \(A\) 为 Noether 整环，\(B\) 为有限型交换 \(A\) 代数，\(M_1,\ldots,M_r\) 为预先给定的有限组有限生成 \(B\) 模。存在同一个 \(0\ne f\in A\)，使 \(B_f\) 与全部 \((M_i)_f\) 都是自由 \(A_f\) 模，因而平坦。
\end{corollary}
\begin{proof}
把有限 \(B\) 模 \(B\oplus\bigoplus_iM_i\) 的每个直和项分别应用泛自由性，取所得非零元素的乘积为 \(f\)。进一步局部化保留自由基，所以各项均自由。自由模平坦已在\cref{b3:def:flat-module}之后直接证明。
\end{proof}

整环 \(A\) 的零素理想属于 \(D(f)\)，且每个非空主开集 \(D(g)\) 都与它相交，因为 \(fg\ne0\)。因此 \(D(f)\) 稠密。这里得到了一个稠密主开集上的平坦性，没有声称整个平坦点集一定由这一个主开集描述。

\subsection{一般纤维与特殊纤维的代数描述}
\label{b3:subsec:1204:04}

对 \(\mathfrak p\in\operatorname{Spec}A\)，代数及模的纤维分别是
\[
B(\mathfrak p)=B\otimes_A\kappa(\mathfrak p),\qquad
M(\mathfrak p)=M\otimes_A\kappa(\mathfrak p).
\]
零素理想处称为一般纤维，使用的是 \(K\)。若 \(M_f\) 已有 \(A_f\) 自由基 \(E\)，则对每个 \(\mathfrak p\in D(f)\)，同一组基的像给出
\(M(\mathfrak p)\cong\kappa(\mathfrak p)^{(E)}\)。因此此开集上的模纤维具有同一基集的大小；当 \(E\) 无限时，不把它误写成某个有限的纤维维数。

\ProseParagraphBreak
例如 \(A=k[t]\)、\(B=A[x]/(tx)\)。在 \(t\ne0\) 处，\(x=0\)，所以 \(B_t\cong A_t\)；在 \(t=0\) 处，纤维却为 \(k[x]\)。作为 \(A\) 模，\(x\ne0\) 而 \(tx=0\)，故 \(B\) 有挠，不平坦。若换成 \(A[x]/(x^2-t)\)，首一除法给出全局基 \(1,x\)，因而处处平坦。即使在 \(t=0\) 处纤维为非约化环 \(k[x]/(x^2)\)，平坦性仍成立；平坦与纤维约化是两种不同的要求。

\subsection{代数有限呈示、模有限展示时的非 Noether 推广}
\label{b3:subsec:1204:05}

\begin{proposition}\label{b3:prop:generic-free-finite-presentation}
设 \(A\) 为任意整环，\(B\) 为有限呈示交换 \(A\) 代数，\(M\) 为有限展示 \(B\) 模。则存在 \(0\ne f\in A\)，使 \(M_f\) 作为 \(A_f\) 模自由；其自由秩允许为零或无限。
\end{proposition}
\begin{proof}
写 \(B=A[x_1,\ldots,x_s]/(h_1,\ldots,h_u)\)，并取有限展示 \(B^v\to B^q\to M\to0\)。把第一映射的有限列提升到 \(P^q\)，再加上 \(h_a e_j\)（\(1\le a\le u,1\le j\le q\)），这些有限列恰好生成 \(P^q\to M\) 的核：任何核元素先模掉 \((h_a)P^q\)，成为 \(B^q\) 中给定关系的组合，提升该组合后差只属于 \((h_a)P^q\)。因此 \(M\) 在 \(P\) 上有限展示。

前述证明在得到有限组 \(v_j\) 后不再使用 \(A\) 的 Noether 性：分式域仍存在，模首项有限性在域上的多项式模中独立成立，且只清除有限个分母。将该证明逐步应用便得结论。
\end{proof}

非 Noether 版本改变的是假设的表达：不能仅凭“有限型代数、有限生成模”自动断言关系有限。有限呈示与有限展示分别保证代数关系和模关系都能由有限数据控制。这项推广不改变自由秩可以无限的结论。
